% Terjemahan appendix1.tex, baris sumber 284--371.
\section{Teorema Gabriel--Popescu}\label{sec:GP}
Bagian ini melanjutkan dan melengkapi \S\ref{sec:Grothendieck-cat}.
Misalkan $\mathcal{A}$ suatu kategori Grothendieck; selanjutnya singkat
$\Hom_{\mathcal{A}}$ menjadi $\Hom$. Notasikan dengan $X^{\oplus I}$ jumlah
langsung di $\mathcal{A}$ dari $I$ salinan objek $X$, dengan $I$ suatu
himpunan kecil.

Teorema Gabriel--Popescu akan menunjukkan bahwa $\mathcal{A}$ selalu dapat direalisasikan sebagai suatu pelokalan reflektif (Catatan \ref{rem:reflexive-localization}) dari kategori modul kanan $\cated{Mod}R$ untuk suatu gelanggang $R$; gelanggang $R$ dan funktor-funktor yang terlibat bergantung pada pembangkit $s$. Di bawah ini kita mengikuti pembuktian L.\ Kuhn \cite{Ku94}, yang pada salah satu langkahnya memerlukan teori klasik funktor turunan; lihat \S\ref{sec:derived-primer}.

Ambil $s \in \Obj(\mathcal{A})$, definisikan $R := \End(s)$, dan tinjau funktor
\begin{align*}
	G: \mathcal{A} & \to \cated{Mod}R \\
	X & \mapsto \Hom(s, X),
\end{align*}
di mana $R$ beraksi dari kanan pada $\Hom(s, X)$ melalui komposisi morfisme. Mudah dilihat bahwa $G$ merupakan funktor aditif di antara kategori-kategori abelian dan mempertahankan semua $\varprojlim$ kecil. Selanjutnya, tulis $\Hom_R := \Hom_{\cated{Mod}R}$ untuk membedakannya dari $\Hom := \Hom_{\mathcal{A}}$.

Untuk setiap $r \in R$, definisikan elemen $\lambda_r: x \mapsto rx$ dari $\End_R(R)$. Perhatikan bahwa $G(s) = R$, sedangkan $\lambda_r: R \to R$ tidak lain adalah citra morfisme $r: s \to s$ di bawah funktor $G$.

\begin{lemma}\label{prop:GP-prep}
	Funktor di atas $G: \mathcal{A} \to \cated{Mod}R$ memiliki adjoin kiri $P: \cated{Mod}R \to \mathcal{A}$. Kounit $\varepsilon: PG \to \identity$ dari pasangan adjoin tersebut membuat $\varepsilon_s: PG(s) = P(R) \to s$ menjadi suatu isomorfisme, dan diagram berikut komutatif:
	\[\begin{tikzcd}[column sep=small]
		\lambda_r \arrow[phantom, r, "\in" description] & \End_R(R) \arrow[r, "P"] & \End(P(R)) \arrow[d, "\sim" sloped] & \varphi \arrow[phantom, l, "\ni" description] \arrow[mapsto, d] \\
		r \arrow[phantom, r, "\in" description] \arrow[mapsto, u] & R \arrow[u, "\sim" sloped] \arrow[equal, r] & \End(s) & \varepsilon_s \varphi \varepsilon_s^{-1} \arrow[phantom, l, "\ni" description] .
	\end{tikzcd}\]
\end{lemma}
\begin{proof}
	Keberadaan adjoin kiri $P$ mengikuti dari teorema funktor adjoin khusus \ref{prop:SAFT} dan Korolari \ref{prop:Grothendieck-cat-wellpowered}. Selanjutnya, untuk setiap $Y \in \Obj(\mathcal{A})$, kita menyatakan bahwa diagram berikut komutatif:
	\[\begin{tikzcd}
		& \Hom(PG(s), Y) \arrow[d, "\sim" sloped] & \\
		\Hom(s, Y) \arrow[ru, "{\varepsilon_s^*}"] \arrow[r, "G"] \arrow[rd, "\identity"'] & \Hom_R(G(s), G(Y)) \arrow[d, "\sim" sloped] & \psi \arrow[phantom, l, "\ni" description, sloped] \arrow[mapsto, d] \\
		& G(Y) & \psi(1_R). \arrow[phantom, l, "\ni" description, sloped]
	\end{tikzcd}\]
	\begin{itemize}
		\item Komutativitas bagian atas merupakan sifat umum pasangan adjoin $(P, G)$: untuk sebarang $f \in \Hom(s, Y)$, morfisme $f\varepsilon_s: PGs \to Y$ berkorespondensi dengan $G(f\varepsilon_s) \circ \eta_{Gs}: Gs \to GY$; lihat \cite[(2.5)]{Li1}. Berdasarkan identitas segitiga untuk pasangan adjoin, morfisme yang terakhir ini sama dengan $(Gf) \circ \left((G\varepsilon_s) \eta_{Gs}\right) = Gf$.
		\item Komutativitas bagian bawah dapat diperiksa langsung dari definisi $G$.
	\end{itemize}
	
	Dengan demikian, $\varepsilon_s^*$ adalah isomorfisme. Selanjutnya, ambil $Y$ sebagai kogenerator injektif yang diberikan oleh Korolari \ref{prop:Grothendieck-cat-cogenerator}; maka terlihat bahwa $\varepsilon_s$ adalah isomorfisme. Terakhir, pernyataan mengenai komutativitas diagram setara dengan mengatakan bahwa
	\[\begin{tikzcd}
		P(R) \arrow[r, "{P(\lambda_r)}"] \arrow[d, "{\varepsilon_s}"'] & P(R) \arrow[d, "{\varepsilon_s}"] \\
		s \arrow[r, "r"'] & s
	\end{tikzcd} \quad (r \in R) \]
	komutatif. Namun, $R = Gs$ dan $\lambda_r = Gr$, sehingga kenaturalan $\varepsilon$ menjamin komutativitas diagram di atas.
\end{proof}

Bagaimana cara mendeskripsikan $P$? Untuk sebarang modul kanan $R$ berupa $M$, terdapat himpunan-himpunan kecil $I$, $J$ dan barisan eksak
\[ R^{\oplus J} \to R^{\oplus I} \to M \to 0; \]
karena $P$ mempertahankan semua $\varinjlim$ kecil, diperoleh $PM \simeq \Coker\left[ s^{\oplus J} \to s^{\oplus I} \right]$.

Pernyataan hasil berikut melibatkan hasil bagi Serre dari suatu kategori abelian; lihat Teorema \ref{prop:Serre-quotient} untuk rinciannya. Hasil bagi Serre merupakan kasus khusus pelokalan.

\begin{theorem}[P.\ Gabriel, N.\ Popescu, L.\ Kuhn]\label{prop:GP}
	Ambil $s$ sebagai pembangkit kategori Grothendieck $\mathcal{A}$. Sifat-sifat berikut berlaku:
	\begin{itemize}
		\item funktor $G = \Hom(s, \cdot): \mathcal{A} \to \cated{Mod}R$ memiliki adjoin kiri eksak $P$;
		\item $G$ penuh dan setia, dan morfisme kounit dari pasangan adjoin memberikan isomorfisme $\varepsilon: PG \rightiso \identity_{\mathcal{A}}$;\footnote{Catatan penerjemah: sumber mencetak $\identity_{\cated{Mod}R}$. Karena $PG$ adalah endofunktor pada $\mathcal{A}$, sasaran kounit yang bertipe benar ialah $\identity_{\mathcal{A}}$.}
		\item $P$ menginduksi ekuivalensi kategori $\cated{Mod}R / \Ker(P) \rightiso \mathcal{A}$.
	\end{itemize}
\end{theorem}
\begin{proof}
	Lema \ref{prop:GP-prep} telah menunjukkan bahwa $G$ memiliki adjoin kiri $P$. Kita menyatakan bahwa jika $u: M \to GX$ adalah monomorfisme dalam $\cated{Mod}R$, maka morfisme yang bersesuaian $v := \varepsilon_X \circ Pu : PM \to X$ (lihat \cite[(2.5)]{Li1}) adalah monomorfisme dalam $\mathcal{A}$. Tuliskan $M$ sebagai $\varinjlim$ terfilter dari submodul-submodul $R$ yang dibangkitkan secara hingga. Karena $P$ mempertahankan $\varinjlim$ dan $\varinjlim$ terfilter dalam $\mathcal{A}$ bersifat eksak, masalahnya tereduksi pada kasus ketika $M$ dibangkitkan secara hingga.
	
	Identifikasikan $s$ dengan $P(R)$ melalui isomorfisme $\varepsilon_s$. Ambil suatu himpunan hingga $I$ beserta $R^{\oplus I} \twoheadrightarrow M$. Ini memberikan $e: s^{\oplus I} \twoheadrightarrow PM$. Untuk menyimpulkan bahwa $v$ merupakan monomorfisme, cukup dibuktikan bahwa $\Ker(ve) = \Ker(e)$. Berdasarkan sifat pembangkit (misalnya, Proposisi \ref{prop:Abel-strong-generator}), masalahnya lebih lanjut tereduksi menjadi menunjukkan bahwa untuk setiap morfisme $f: s \to s^{\oplus I}$ yang memenuhi $vef = 0$, berlaku $ef = 0$. Buat dua pengamatan:
	\begin{itemize}
		\item Menurut konstruksinya, $e$ merupakan citra suatu morfisme di bawah $P$. Karena $I$ hingga, tuliskan $f$ sebagai suatu elemen dari $R^{\oplus I}$ dan periksa komponennya satu per satu; diagram komutatif dalam Lema \ref{prop:GP-prep} menjamin bahwa $f$ juga merupakan citra di bawah $P$;
		\item Selanjutnya, jika morfisme $\psi: N \to M$ dalam $\cated{Mod}R$ memenuhi $v \circ P\psi = 0$, maka $\psi = 0$. Ini mengikuti dari monomorfisitas $u$ dan diagram komutatif dalam $\cate{Ab}$
		\[\begin{tikzcd}[column sep=small]
			v \arrow[phantom, r, "\in" description]& \Hom(PM, X) \arrow[r, "\sim"] \arrow[d, "{(P\psi)^*}"'] & \Hom_R(M, GX) \arrow[d, "{\psi^*}"] & u \arrow[phantom, l, "\ni" description] \\
			& \Hom(PN, X) \arrow[r, "\sim"'] & \Hom_R(N, GX) . &
		\end{tikzcd}\]
	\end{itemize}
	Dengan demikian, $ef: s \to PM$ selalu merupakan citra suatu morfisme di bawah $P$, sedangkan $vef = 0$ menyiratkan $ef=0$. Pernyataan tersebut terbukti.
	
	Morfisme kounit $\varepsilon$ dari pasangan adjoin adalah isomorfisme jika dan hanya jika $G$ penuh dan setia; fakta umum ini merupakan latihan dalam \CHref{sec:categories}. Sekarang kita buktikan bahwa $\varepsilon$ adalah isomorfisme. Ambil $X \in \Obj(\mathcal{A})$. Karena $\varepsilon_X: PG(X) \hookrightarrow X$ berkorespondensi dengan $\identity_{GX}$, dengan mengambil $u = \identity_{GX}$ pada langkah sebelumnya diketahui bahwa $\varepsilon_X$ merupakan monomorfisme. Untuk membuktikan bahwa ia epimorfisme, Proposisi \ref{prop:Abel-strong-generator} mereduksi masalahnya menjadi membuktikan bahwa setiap $\alpha: s \to X$ berfaktor melalui $\varepsilon_X: PG(X) \hookrightarrow X$, dan diagram berikut memberikan faktorisasi yang diminta:
	\[\begin{tikzcd}
		PG(s) \arrow[r, "\sim"', "\varepsilon_s"] \arrow[d, "{PG(\alpha)}"'] & s \arrow[d, "\alpha"] \\
		PG(X) \arrow[r, "{\varepsilon_X}"'] & X
	\end{tikzcd}\]
	
	Selanjutnya, kita buktikan bahwa $P$ eksak. Telah diketahui bahwa $P$ eksak kanan, sehingga cukup dibuktikan bahwa funktor turunan kiri $\mathrm{L}_1 P = 0$. Untuk sebarang modul kanan $R$ berupa $M$, ambil barisan eksak pendek
	\[ 0 \to K \xrightarrow{u} F \to M \to 0, \quad F: \text{modul bebas}. \]
	Berdasarkan teknik pergeseran dimensi (Proposisi \ref{prop:dimension-shifting}), diperoleh
	\[ \mathrm{L}_1 P (M) \simeq \Ker[Pu: P(K) \to P(F)]. \]
	Dengan demikian, masalahnya tereduksi menjadi membuktikan bahwa $Pu$ merupakan monomorfisme. Tuliskan $F$ sebagai $\varinjlim$ terfilter dari submodul-submodul bebas $F'$ berperingkat hingga (atau sebagai gabungan menaik), dan secara bersesuaian $K = \varinjlim_{F'} (F' \cap K)$. Karena $P$ mempertahankan $\varinjlim$ dan $\varinjlim$ terfilter dalam $\mathcal{A}$ bersifat eksak, masalahnya tereduksi pada kasus $F = R^{\oplus I} \simeq G(s^{\oplus I})$, dengan $I$ suatu himpunan hingga. Akan tetapi, pada awal pembuktian telah ditunjukkan bahwa jika $u: K \to G(s^{\oplus I})$ merupakan monomorfisme, maka $v = \varepsilon_{s^{\oplus I}} \circ Pu : P(K) \to s^{\oplus I}$ merupakan monomorfisme; akibatnya, $Pu$ juga merupakan monomorfisme. Keeksakan pun terbukti.
	
	Terakhir, dengan menerapkan sifat universal yang diberikan oleh Teorema \ref{prop:Serre-quotient} pada $P$, kita dapat memfaktorkan $P$ melalui suatu funktor aditif setia $\overline{P}: \cated{Mod}R / \Ker(P) \to \mathcal{A}$. Tuliskan $\overline{G}$ untuk komposisi $\mathcal{A} \xrightarrow{G} \cated{Mod}R \to \cated{Mod}R / \Ker(P)$. Maka $\overline{P} \circ \overline{G} \simeq P \circ G \simeq \identity_{\mathcal{A}}$. Dari sini diketahui bahwa $\overline{P}$ juga penuh, setia, dan surjektif secara esensial, dengan $\overline{G}$ sebagai kuasi-inversnya. Ini membuktikan pernyataan tersebut.
\end{proof}

\begin{remark}
	Berdasarkan Teorema Gabriel--Popescu \ref{prop:GP}, dapat disimpulkan bahwa kategori Grothendieck merupakan kategori presentabel dalam pengertian Definisi \ref{def:presentable-cat}. Ini karena $\cated{Mod}R \simeq R^{\opp}\dcate{Mod}$ diketahui presentabel; dengan menerapkan \cite[Korolari 1.40]{AR94}, sifat ini dapat ``direfleksikan'' ke $\mathcal{A}$.
\end{remark}
